% !TEX TS-program = pdflatex
% !TEX encoding = UTF-8 Unicode

% This file is a template using the "beamer" package to create slides for a talk or presentation
% - Giving a talk on some subject.
% - The talk is between 15min and 45min long.
% - Style is ornate.

% MODIFIED by Jonathan Kew, 2008-07-06
% The header comments and encoding in this file were modified for inclusion with TeXworks.
% The content is otherwise unchanged from the original distributed with the beamer package.

\documentclass{beamer}
\setbeamercovered{invisible}



% Copyright 2004 by Till Tantau <tantau@users.sourceforge.net>.
%
% In principle, this file can be redistributed and/or modified under
% the terms of the GNU Public License, version 2.
%
% However, this file is supposed to be a template to be modified
% for your own needs. For this reason, if you use this file as a
% template and not specifically distribute it as part of a another
% package/program, I grant the extra permission to freely copy and
% modify this file as you see fit and even to delete this copyright
% notice. 


\mode<presentation>
{
  \usetheme{Warsaw}
  % or ...

  % or whatever (possibly just delete it)
}


\usepackage[english]{babel}
% or whatever

\usepackage[utf8]{inputenc}
\usepackage{qrcode}
\usepackage{hyperref}
% or whatever

\usepackage{times}
\usepackage[T1]{fontenc}
% Or whatever. Note that the encoding and the font should match. If T1
% does not look nice, try deleting the line with the fontenc.

%%% MATH RELATED

\title{MATH 180 Strategies of Problem Solving} 
\subtitle
{} % (optional)

\author[W.R. Casper] % (optional, use only with lots of authors)
{W.R. Casper}
% - Use the \inst{?} command only if the authors have different
%   affiliation.

\institute[California State University Fullerton] % (optional, but mostly needed)
{
  Department of Mathematics\\
  California State University Fullerton}
% - Use the \inst command only if there are several affiliations.
% - Keep it simple, no one is interested in your street address.

\subject{Talks}
% This is only inserted into the PDF information catalog. Can be left
% out. 



% If you have a file called "university-logo-filename.xxx", where xxx
% is a graphic format that can be processed by latex or pdflatex,
% resp., then you can add a logo as follows:

% \pgfdeclareimage[height=0.5cm]{university-logo}{university-logo-filename}
% \logo{\pgfuseimage{university-logo}}



% Delete this, if you do not want the table of contents to pop up at
% the beginning of each subsection:
\AtBeginSubsection[]
{
  \begin{frame}<beamer>{Outline}
    \tableofcontents[currentsection,currentsubsection]
  \end{frame}
}


% If you wish to uncover everything in a step-wise fashion, uncomment
% the following command: 

%\beamerdefaultoverlayspecification{<+->}


\begin{document}

\begin{frame}
  \titlepage
\end{frame}

\begin{frame}{Outline}
  \tableofcontents
  % You might wish to add the option [pausesections]
\end{frame}

% Since this a solution template for a generic talk, very little can
% be said about how it should be structured. However, the talk length
% of between 15min and 45min and the theme suggest that you stick to
% the following rules:  

% - Exactly two or three sections (other than the summary).
% - At *most* three subsections per section.
% - Talk about 30s to 2min per frame. So there should be between about
%   15 and 30 frames, all told.

\section{SoPS Lecture 1}
\subsection{Welcome}

\begin{frame}{What is this course about?}
\pause
  \begin{block}{Learn what mathematicians do}
    \begin{itemize}
\pause
      \item play $\rightarrow$ discover patterns
\pause
      \item abstract $\rightarrow$ express patterns mathematically
\pause
      \item prove $\rightarrow$ justify why the pattern is true
    \end{itemize}
  \end{block}
\pause
  \begin{block}{Glimpse advanced mathematics}
    \begin{itemize}
\pause
      \item logic, set theory, cardinality
\pause
      \item probability, combinatorics, graph theory
    \end{itemize}
  \end{block}
\pause
  \begin{block}{Make friends}
    \begin{itemize}
\pause
      \item math is {\color{red}social}
\pause
      \item you'll see each other \emph{a lot}
    \end{itemize}
  \end{block}
\end{frame}

\begin{frame}{Course expectations}
\pause
  \begin{block}{Grading}
    \begin{itemize}
\pause
      \item pre-class/post-class assignments $\rightarrow 15\%$
\pause
      \item in-class assignments $\rightarrow 25\%$ 
\pause
      \item homework $\rightarrow 25\%$ 
\pause
      \item seminars / research $\rightarrow 15\%$ 
\pause
      \item final exam $\rightarrow 20\%$ 
    \end{itemize}
  \end{block}
\end{frame}

\begin{frame}{Course expectations}
\pause
  \begin{block}{Daily Expectations}
    \begin{itemize}
\pause
      \item complete pre-class/post-class assignments on time
\pause
      \item attend every class 
\pause
      \item participate \textbf{actively}
\pause
      \item solicit participation from group members
\pause
      \item do your homework (or else)
    \end{itemize}
  \end{block}
\end{frame}

\begin{frame}{Course expectations}
\pause
  \begin{block}{Seminar Participation}
    \begin{itemize}
\pause
      \item must attend at least three
\pause
      \item write a one-page report
\pause
      \item your time at CSUF is what you make it!
    \end{itemize}
  \end{block}
\pause
  \begin{block}{Ongoing Seminars}
    \begin{itemize}
\pause
      \item pure math seminar
\pause
      \item applied math seminar
\pause
      \item math education seminar
\pause
      \item problem solving seminar
\pause
      \item other events (workshops, colloquia, etc)
    \end{itemize}
  \end{block}
\end{frame}

\begin{frame}{Course expectations}
\pause
  \begin{block}{Research project}
    \begin{itemize}
\pause
      \item what math research is
\pause
      \item do original research
\pause
      \item write a final report
\pause
      \item present results
    \end{itemize}
  \end{block}
\end{frame}

\subsection{New game}
\begin{frame}{Pinching Pennies}
\pause
  \begin{block}{Rules}
    \begin{itemize}
\pause
      \item arrange pennies in two piles
\pause
      \item turn: take one or more pennies \emph{but}
\pause
      \item take only from one pile per turn
\pause
      \item must take at least one penny
\pause
      \item person who takes last penny wins
    \end{itemize}
  \end{block}
\pause
  \begin{block}{Activity (20 mins)}
    \begin{itemize}
\pause
      \item split into pairs
\pause
      \item play 4-5 rounds
\pause
      \item {\color{red}record data:} initial board, who won? 
    \end{itemize}
  \end{block}
\end{frame}

\begin{frame}{Pinching Pennies: Debrief}
\pause
\textbf{Let's combine the data ...}
\pause
  \begin{block}{Question 1}
     is it better to be {\color{blue}Player 1} or {\color{red}Player 2} ?
  \end{block}
\pause
  \begin{block}{Question 2}
     if the piles are equal, then who should win?
  \end{block}
\pause
\begin{center}
\textbf{Fight me!}
\end{center}
\end{frame}

\begin{frame}{Vocabulary}
\pause
  \begin{block}{Optimal Play}
\pause
    A player is said to \textbf{play optimally} or to use an \textbf{optimal strategy} if their moves put them in the best possible final position, regardless of the moves of other players.
  \end{block}
\pause
  \begin{block}{Winning/Losing Position}
\pause
    A \textbf{winning position} for a player in a game is a state of the game where if that player {plays optimally}, they are guaranteed to win.  A \textbf{losing position} for a player is a game state where one of their opponents has a winning position.
  \end{block}
\end{frame}

\subsection{New game plus}
\begin{frame}{Pinching Pennies -- revisited}
\pause
  \begin{block}{Activity (20 minutes)}
    \begin{itemize}
\pause
      \item arrange pennies in \textbf{three piles}
\pause
      \item which starts are losing positions for {\color{blue}Player 1}?
    \end{itemize}
  \end{block}
\pause
  \begin{block}{Submit your conjecture:}
    \begin{center}
      \qrcode[hyperlink,height=2cm]{https://padlet.com/williamrileycasper/pinching-pennies-losing-positions-h5xyswjj3ku3iqof}
    \end{center}
  \end{block}
\end{frame}

\begin{frame}{Break time!}
  \begin{block}{Activity (5 minutes)}
    \begin{itemize}
      \item get up, stretch get a cup of coffee
      \item chat up your neighbor, introduce yourself
    \end{itemize}
  \end{block}
\end{frame}

\begin{frame}{Pinching Pennies -- two equal piles}
\pause
  \begin{block}{Question 1:}
    Two piles with one penny each, who wins?
  \end{block}
\pause
  \begin{block}{Question 2:}
    Two piles with two pennies each, who wins?
  \end{block}
\pause
  \begin{block}{Question 3:}
    Two piles with three pennies each, who wins?
  \end{block}
\pause
  \begin{block}{Copycat strategy}
     {\color{red}Player 2} copies what {\color{blue}Player 1} does to the opposite pile
  \end{block}
\end{frame}

\begin{frame}{Our first proof}
\pause
  \begin{block}{Theorem}
    Two equal piles of pennies is a losing position for {\color{blue}Player 1}.
  \end{block}
\pause
  \begin{block}{Proof}
\pause
    Start with $n$ pennies in each pile.\\
\pause
    If {\color{blue}Player 1} takes a whole pile, {\color{red}Player 2} wins.\\
\pause
    Otherwise {\color{blue}Player 1} must reduce a pile to $m$ pennies with $m<n$.\\
\pause
    Then {\color{red}Player 2} reduces the other pile to $m$ pennies.\\
\pause
    Now we're back where we started, but with less pennies!\\
\pause
    Repeat until {\color{blue}Player 1} is forced to take a whole pile.
  \end{block}
\end{frame}

\subsection{Finishing up}
\begin{frame}{Question of the Day}
  \begin{block}{Puzzle}
    You have 15 pennies. On each turn, a player removes 1, 2, or 3 coins.  The player who takes the last coin wins. 
    Is the start a winning position for {\color{blue}Player 1} or {\color{red}Player 2}?
  \end{block}
\end{frame}

\begin{frame}{Final thoughts}
  \begin{block}{Reflection}
    \begin{itemize}
      \item What did you notice about how mathematicians work compared to how you expected?
      \item What differences do you see between experimentation (play) and proof?
      \item How is this class different from the math you’ve done before?  What do you think is the key to your success in this class?
    \end{itemize}
  \end{block}
\begin{center}
Submit your responses by \textbf{tonight!}
\end{center}
\end{frame}


\end{document}


